% =============================================================================
% Odd Ratio Linear Fit - LaTeX Description for Papers
% =============================================================================
% This file contains a description of the odd ratio linear fit method
% suitable for copy-pasting into the Methods section of a scientific paper.
%
% Reference: Artigau et al. (2022), AJ, 164, 84, Appendix A
% =============================================================================

% -----------------------------------------------------------------------------
% OPTION 1: Compact version (~150 words)
% Use this for papers with strict space constraints
% -----------------------------------------------------------------------------

% \subsection{Robust Linear Fitting}
%
% We performed linear regression using the odd ratio mixture model method
% \citep{Artigau2022}, which is robust to outliers. This approach models each
% data point as coming from either a Gaussian distribution centered on the
% true model (valid measurements) or a uniform background distribution
% (outliers). For each point with residual $r_i$ and uncertainty $\sigma_i$,
% the probability of being valid is computed as:
% \begin{equation}
%     p_i = 1 - \frac{f_0}{(1-f_0) \exp\left(-r_i^2/2\sigma_i^2\right) + f_0}
% \end{equation}
% where $f_0 = 2 \times 10^{-4}$ is the prior probability that any point is
% an outlier. The fit is performed iteratively: weights $w_i = p_i/\sigma_i^2$
% are computed from the current fit, then weighted least squares is applied
% with these weights, repeating until convergence (typically 3--5 iterations).
% This method avoids arbitrary sigma-clipping thresholds while preserving
% proper error propagation.


% -----------------------------------------------------------------------------
% OPTION 2: Standard version (~300 words)
% Recommended for most papers
% -----------------------------------------------------------------------------

\subsection{Robust Linear Fitting with Mixture Model}
\label{sec:odd_ratio_fit}

We employed the odd ratio mixture model method (Appendix A of \citealt{Artigau2022}) for robust
linear regression. Unlike traditional sigma-clipping approaches that use hard
thresholds to reject outliers, this method uses a probabilistic model that
naturally down-weights anomalous points while preserving proper statistical
uncertainties.

The algorithm models each measurement as arising from a two-component mixture:
valid measurements are drawn from a Gaussian distribution centered on the true
model, while outliers come from a uniform background distribution. Given a
linear model $y = a + bx$ and residual $r_i = y_i - (a + bx_i)$ for point $i$
with measurement uncertainty $\sigma_i$, the posterior probability that the
point is a valid measurement is:
\begin{equation}
    p_{\mathrm{valid},i} = 1 - \frac{f_0}{(1-f_0) \exp\left(-r_i^2 / 2\sigma_i^2\right) + f_0}
    \label{eq:odd_ratio}
\end{equation}
where $f_0$ is the prior probability that any given point is an outlier. We
adopted $f_0 = 2 \times 10^{-4}$ following \citet{Artigau2022}.

The fitting procedure is iterative:
\begin{enumerate}
    \item Initialize with standard weighted least squares using weights $w_i = 1/\sigma_i^2$
    \item Compute residuals from the current best-fit model
    \item Calculate $p_{\mathrm{valid},i}$ for each point using Equation~\ref{eq:odd_ratio}
    \item Update weights to $w_i = p_{\mathrm{valid},i} / \sigma_i^2$
    \item Re-fit using weighted least squares with the updated weights
    \item Repeat steps 2--5 until convergence
\end{enumerate}

The algorithm typically converges in 3--5 iterations. Parameter uncertainties
are computed from the weighted covariance matrix, where weights reflect the
probability of each point being valid. This approach provides several
advantages over sigma-clipping: it avoids arbitrary threshold choices, allows
points to be partially down-weighted rather than completely rejected, ensures
stable convergence (binary inclusion/rejection can cause oscillations where
points flip in and out of the fit), and maintains proper error propagation
throughout.

The method is implemented in the \texttt{odd\_ratio\_fits} Python package\footnote{%
    \url{https://github.com/eartigau/odd_ratio_fits}}.


% -----------------------------------------------------------------------------
% OPTION 3: Extended version with mathematical details (~500 words)
% Use when space permits or for methodology-focused papers
% -----------------------------------------------------------------------------

% \subsection{Robust Linear Fitting with Odd Ratio Mixture Model}
% \label{sec:odd_ratio_extended}
%
% We employed a robust linear fitting method based on the odd ratio mixture
% model introduced in Appendix A of \citet{Artigau2022}. This approach treats outlier
% rejection as a Bayesian inference problem rather than applying hard
% thresholds, providing a principled framework that maintains proper
% statistical uncertainties.
%
% \subsubsection{Mixture Model Formulation}
%
% Consider fitting a linear model $y = a + bx$ to data $\{x_i, y_i, \sigma_i\}$
% where $\sigma_i$ is the measurement uncertainty. We model each data point
% as arising from a two-component mixture:
% \begin{equation}
%     P(y_i | x_i, a, b) = (1-f_0) \cdot \mathcal{N}(y_i | a+bx_i, \sigma_i) + f_0 \cdot P_{\mathrm{outlier}}
% \end{equation}
% where $\mathcal{N}$ denotes a Gaussian distribution, $f_0$ is the prior
% probability that any point is an outlier, and $P_{\mathrm{outlier}}$ represents
% the outlier distribution (assumed to be uniform and broad).
%
% For a point with residual $r_i = y_i - (a + bx_i)$, the posterior probability
% that it is a valid measurement (not an outlier) is:
% \begin{equation}
%     p_{\mathrm{valid},i} = \frac{(1-f_0) \exp\left(-r_i^2/2\sigma_i^2\right)}
%         {(1-f_0) \exp\left(-r_i^2/2\sigma_i^2\right) + f_0}
%     \label{eq:p_valid}
% \end{equation}
%
% Following \citet{Artigau2022}, we adopt $f_0 = 2 \times 10^{-4}$. This value
% implies that roughly 1 in 5000 points is expected to be an outlier \emph{a priori},
% which is conservative for most astronomical datasets. Points deviating by
% $\gtrsim 3.5\sigma$ are significantly down-weighted.
%
% \subsubsection{Iterative Fitting Algorithm}
%
% The fitting proceeds iteratively:
% \begin{enumerate}
%     \item \textbf{Initialization}: Perform standard weighted least squares
%           with $w_i = 1/\sigma_i^2$ to obtain initial estimates $\hat{a}$, $\hat{b}$.
%     \item \textbf{E-step}: Compute $p_{\mathrm{valid},i}$ for each point
%           using Equation~\ref{eq:p_valid} with current parameter estimates.
%     \item \textbf{M-step}: Update the model parameters using weighted least
%           squares with weights $w_i = p_{\mathrm{valid},i}/\sigma_i^2$:
%           \begin{align}
%               \hat{b} &= \frac{\sum w_i (x_i-\bar{x})(y_i-\bar{y})}{\sum w_i (x_i-\bar{x})^2} \\
%               \hat{a} &= \bar{y} - \hat{b}\bar{x}
%           \end{align}
%           where $\bar{x} = \sum w_i x_i / \sum w_i$ and similarly for $\bar{y}$.
%     \item \textbf{Convergence check}: If the relative change in parameters
%           is below a threshold (default: $10^{-2}$ times the parameter uncertainty),
%           stop; otherwise return to step 2.
% \end{enumerate}
%
% The algorithm typically converges in 3--5 iterations. Parameter uncertainties
% are derived from the weighted covariance matrix:
% \begin{align}
%     \sigma_a^2 &= \frac{\sum w_i x_i^2}{\Delta} \\
%     \sigma_b^2 &= \frac{\sum w_i}{\Delta}
% \end{align}
% where $\Delta = (\sum w_i)(\sum w_i x_i^2) - (\sum w_i x_i)^2$.
%
% This method offers key advantages over sigma-clipping: (1) no arbitrary
% threshold selection, (2) smooth weighting allows partial down-weighting
% of marginal points, and (3) proper error propagation accounting for
% effective sample size reduction.


% -----------------------------------------------------------------------------
% Bibliography entry
% -----------------------------------------------------------------------------
% Add this to your .bib file:
%
% @ARTICLE{Artigau2022,
%        author = {{Artigau}, {\'E}tienne and {Cadieux}, Charles and {Cook}, Neil J. and
%          {Doyon}, Ren{\'e} and {Vandal}, Thomas and {Donati}, Jean-Fran{\c{c}}ois and
%          {Moutou}, Claire and {Delfosse}, Xavier and {Fouqu{\'e}}, Pascal and
%          {Collier Cameron}, Andrew and others},
%         title = "{Line-by-line Velocity Measurements: an Outlier-resistant Method for Precision Velocimetry}",
%       journal = {\aj},
%          year = 2022,
%         month = sep,
%        volume = {164},
%        number = {3},
%           eid = {84},
%         pages = {84},
%           doi = {10.3847/1538-3881/ac7f2c},
% archivePrefix = {arXiv},
%        eprint = {2207.13524},
% }
